\section{An exactly normalized pure energy coefficient}
\label{sec:pure}

The scale argument needs a coefficient that is preserved exactly when a
local observable is integrated through a pure annulus.  The approximate
coefficients in the companion transfer theorem identify the relevant
one-dimensional direction, but do not by themselves fix its normalization
through infinitely many scales.  We fix that normalization by the known
energy covariance.  We then prove two estimates needed below: cancellation
of the energy coefficient of a product of separated centered observables,
and a transfer bound for an arbitrarily large positive local weight.

Let $\mu_0$ be the full-plane nearest-neighbor Ising state at
$K_0=\frac12\log(1+\sqrt2)$, and write $\E_0$ for its expectation.
Set $B_s(x)=x+[-s,s]^2\cap\Z^2$ and
$\partial B_s(x)=B_s(x)\setminus B_{s-1}(x)$, initially for integer
$s\geq1$, and abbreviate $B_s=B_s(0)$.  Observables in this section are
real-valued.  A local observable is \emph{even} if it is unchanged by
simultaneously reversing all spins.  It may also depend on independent
auxiliary variables attached to its support.  These variables are integrated
before a pure transfer; spin reversal leaves them unchanged.  All norms in
this section are supremum norms, and all centered observables have mean zero
under $\E_0$, including this auxiliary integration.

For an observable supported strictly inside $B_R(x)$ define
\[
 P_{R,x}f=\E_0[f\mid\sigma_{\partial B_R(x)}]-\E_0f,
 \qquad P_R=P_{R,0}.
\]
We use $P_{R,x}^{\mathrm{unc}}$ for the same conditional expectation without
subtracting the mean.  The spatial Markov property gives exact composition
for nested squares.  It also permits simultaneous integration inside
disjoint squares while preserving every correlation with their exterior.
Integer roundings of fixed radius ratios will always be understood.

Critical duality gives $\E_0\sigma_z\sigma_{z'}=1/\sqrt2$ on every
nearest-neighbor edge.  Indeed, differentiating the finite planar
partition-function duality once and passing to the unique critical
full-plane state gives
$\E_0\sigma_z\sigma_{z'}=\coth(2K_0)-\E_0\sigma_z\sigma_{z'}$,
and $\coth(2K_0)=\sqrt2$ \cite{KW}.
For an unordered nearest-neighbor edge $e=\{z,z'\}$ put
\[
 b_e=\sigma_z\sigma_{z'}-\frac1{\sqrt2},\qquad
 B_x=\frac14\sum_{e\ni x}b_e,\qquad B=B_0,
 \qquad \kappa=\frac1{\pi^2}.
\]
The symbol $B_x$ denotes this observable, whereas $B_s(x)$ denotes a square.

\subsection{The companion inputs and their nonsymmetric consequence}

We use the following statements of the companion at the scopes specified
here.  These are pure-model statements throughout.

First, the Proposition ``Buffered and small-patch comparisons'' in
\cite[Section~2]{S} bounds the density of the spin law on one side of a
fixed proportional unpinned collar, under any data on the other side,
between $K^{-1}$ and $K$ relative to the full-plane marginal.  The constant
depends only on the collar ratios.  The assertion holds with the two
terminal sets interchanged and for mixtures of boundary data.  In particular,
conditional expectation through such a collar maps $L^1(\mu_0)$ to bounded
functions with norm at most $K$ times the $L^1$ norm.

Second, the Proposition ``Positive energy covariance'' in
\cite[Section~2]{S} gives, for either orientation of each bond,
\begin{equation}
 \Cov_0(b_e,b_{e'+x})=\frac{\kappa+o(1)}{|x|^2}
 \quad (|x|\longrightarrow\infty).
 \label{pure:energy-input}
\end{equation}
The same formula holds for $B,B_x$.  Fixed displacements of the bonds do
not change its coefficient.

Third, we use the Lemma ``Bounded replacement across a gap'', the radial
Hilbert-space construction, and the Proposition ``Radial spectral gap''
of \cite[Section~3]{S}.  Bounded replacement concerns arbitrary bounded
lattice observables on finitely many positively separated square or band
supports.  Their pairings can be approximated simultaneously by bounded
continuum spin-field tests, with the same supremum bounds.  The menus of
spatial smears are fixed before the mesh tends to zero.  The radial
dilation semigroup has a positive spectral representation on the Hilbert
space generated by spin-field tests supported in the unit disk, and its
Hilbert norm is bounded by ordinary $L^2$.  We denote its inner product
by $(\cdot,\cdot)_{\mathrm{rad}}$ and its inward-dilation semigroup by
$D_t$, with radius ratio $e^{-t}$.  On even quarter-turn-invariant tests its spectrum below $3$
consists of the constant state at $0$ and a single state at $1$.
The full-plane spin-correlation formula and moment bounds used below are
those in the Proposition ``Spin normalization and moment bounds''
of \cite[Section~2]{S}.  Write $S(z_1,\ldots,z_n)$ for the continuum
$n$-point spin correlation in its normalization, so that
$S(z,z')=|z-z'|^{-1/4}$.  A smooth lattice smear at scale $s$ is
\[
 \Phi_s(\psi)=a_\sigma^{-1}s^{-15/8}
                  \sum_{z\in\Z^2}\psi(z/s)\sigma_z,
\]
where $a_\sigma>0$ is the lattice normalization in that proposition.
Every fixed finite collection of these smears has convergent moments of
all orders; the bounds are uniform in $s$.

Here is the consequence of these inputs that also treats observables without
spatial symmetry.  If $\mathcal A_s$ denotes quarter-turn averaging about the
center of $B_s$, the subspace $\ker\mathcal A_s$ means that the average of
the observable, rather than its individual values, vanishes.

\begin{lemma}[Fixed-ratio pure transfer]
\label{pure:fixed-ratio}
There is a constant $C$ with the following property.  Fix finitely many
sufficiently large integer ratios $d$.  A common measurement ratio $H$ and
then a threshold $s_*$ can be chosen so that, for every integer $s\geq s_*$,
there are linear functionals $\ell_s$ on centered even
quarter-turn-invariant observables in $B_s$, of norm at most $C$, and
centered right inverses $e_s$ with $\ell_s(e_s)=1$ and $\|e_s\|_\infty\leq C$.
For each chosen ratio, with $T_{s,d}=P_{ds}$,
\begin{align}
 \|T_{s,d}f\|_\infty
   &\leq C d^{-1}|\ell_s(f)|+C d^{-3}\|f\|_\infty,
       \label{pure:ratio-main}\\
 |\ell_{ds}(T_{s,d}f)-d^{-1}\ell_s(f)|
   &\leq C d^{-3}\|f\|_\infty.
       \label{pure:ratio-coefficient}
\end{align}
If $f$ is even and $\mathcal A_s f=0$, then
\begin{equation}
 \|T_{s,d}f\|_\infty\leq C d^{-2}\|f\|_\infty.
 \label{pure:ratio-nonscalar}
\end{equation}
The constant $C$ is independent of the chosen large ratios.
\end{lemma}

\begin{proof}
For the invariant subspace use the displayed upper-limit estimates in
\cite[Section~3, ``Compatible lattice coefficients and arbitrary boundary
data'']{S}, dividing the companion's density-normalized operator by $d^2$.
The measured coefficient is
\[
 \ell_s(f)=\frac H g\Cov_0(f,W_{Hs}),
\]
where $W$ is the fixed exterior annular test there and $g\ne0$ is its
radial degree-one coefficient.  The remainder in the transfer norm is
$C d^{-3}+C/(dH^2)$, and that in coefficient compatibility is
$C/(dH^2)$.  Choose $H$ large enough for all the prescribed ratios, then
the replacement accuracies, and finally $s_*$ to absorb the uniform
upper-limit errors into $C d^{-3}$.  Center the companion's fixed interior
reference test under $\E_0$ and divide it by its measured coefficient to
obtain the centered bounded right inverse.  Centering changes neither
that coefficient nor the transfer estimates.  The
radial spectral gap requires only quarter turns, so reflection averaging
in this application is unnecessary.

For completeness, the same method proves \eqref{pure:ratio-nonscalar}.
Before the angular averaging in the proof of the companion's radial gap,
the polynomial expansion has the form
\[
 (F,D_{-\log r}G)_{\mathrm{rad}}
   =\E\overline F\,\E G
      +r\,\overline{c(F)}c(G)+O_{F,G}(r^2).
\]
Both displayed coefficients are rotation scalars.  Before smearing, the
second factor is
$c(\prod_i\sigma(z_i))=S(z_1,\ldots,z_{2m})A_2(z)/A_0(z)$, where
\[
 A_q(z)=\sum_{\substack{u_i\in\{-1,1\}\\\sum_i u_i=q}}
            \prod_{i<j}|z_i-z_j|^{u_i u_j/2}.
\]
All distances in this expression are rotation invariant.  Thus both vanish
on the even zero-quarter-turn-average subspace.  A polynomial self-pairing
there is the integral of $r^a$ against a positive spectral measure and is
$O_F(r^2)$.  Positive mass in any compact subinterval of $[0,2)$ would
contradict this bound as $r\downarrow0$.  Polynomial density therefore
places the whole spectrum of this subspace in $[2,\infty)$, giving the
uniform radial bound $r^2\|F\|_2\|G\|_2$.

Apply bounded replacement to the observable and to the buffered boundary
density used in the proof of uniform transfer.  Project the first proxy
with $I-\mathcal A$; this changes its norm and the replacement error by at
most a factor two.  The radial bound just proved then gives
$C d^{-2}$, uniformly over the boundary assignment.  The same finite-list
choice of accuracies and threshold makes the lattice inequality uniform
in $s$ and in the observable.  This proves the lemma without a spatial
symmetry assumption on the original even observable.
\end{proof}

\subsection{Exact normalization by the energy covariance}

Fix
\[
 u=\frac95,\qquad w=\frac{12}5.
\]
We first extract an exact invariant hyperplane from the approximate
coefficient and then calibrate its one-dimensional complement.  The latter
step is what prevents small coefficient errors from accumulating over
infinitely many scales.

\begin{lemma}[The stable hyperplane]
\label{pure:hyperplane}
Fix a sufficiently large ratio $d$ and a geometric sequence $s_j=s_*d^j$.
On centered even quarter-turn-invariant observables in $B_{s_j}$ there
are uniformly bounded linear functionals $\beta_j$ such that
\[
 |\beta_j(f)-\ell_{s_j}(f)|\leq C d^{-2}\|f\|_\infty,
 \qquad
 \beta_{j+1}(P_{s_{j+1}}f)=m_j\beta_j(f),
\]
where $m_j=d^{-1}(1+O(d^{-2}))>0$.
An observable in $\ker\beta_j$ satisfies
\begin{equation}
 \|P_{s_{j+n}}f\|_\infty
      \leq C d^{-wn}\|f\|_\infty\qquad(n\geq1).
 \label{pure:stable-decay}
\end{equation}
Every observable outside this kernel eventually enters an invariant cone
\begin{equation}
 \|f\|_\infty\leq M|\ell_{s_j}(f)|,
 \label{pure:cone}
\end{equation}
where $M$ is independent of $j$ and of sufficiently large $d$.
\end{lemma}

\begin{proof}
Start with a bounded right inverse $v_0$ of $\ell_{s_*}$ and propagate it
by
\[
 m_j=\ell_{s_{j+1}}(P_{s_{j+1}}v_j),\qquad
 v_{j+1}=m_j^{-1}P_{s_{j+1}}v_j.
\]
Equations~\eqref{pure:ratio-main}--\eqref{pure:ratio-coefficient} show by
induction that $\|v_j\|_\infty$ stays bounded and
$m_j=d^{-1}(1+O(d^{-2}))$.  Indeed, choose a fixed bound larger than four
times the constant in \eqref{pure:ratio-main}, then choose $d$ so large
that the $d^{-2}$ errors preserve it.

Write $f_j=a_jv_j+k_j$, where $\ell_{s_j}(k_j)=0$.  Transfer has the
triangular form
\[
 a_{j+1}=m_ja_j+t_jk_j,\qquad k_{j+1}=K_jk_j,
 \qquad \|t_j\|+\|K_j\|\leq C d^{-3}.
\]
The last estimate follows by subtracting
$v_{j+1}\ell_{s_{j+1}}(P_{s_{j+1}}k_j)$ from the transferred kernel
vector.  Put $M_{j,n}=\prod_{h=j}^{n-1}m_h$, with $M_{j,j}=1$, and
define
\[
 \beta_j(a v_j+k)
  =a+\sum_{n=j}^{\infty}
       \frac{t_n K_{n-1}\cdots K_jk}{M_{j,n+1}}.
\]
For $n=j$ the product of $K$'s is the identity.  The series converges in
operator norm because the ratio of its successive bounds is $O(d^{-2})$;
its correction to $a$ has norm $O(d^{-2})\|k\|_\infty$.
Shifting the series proves exact compatibility with multiplier $m_j$.

If $\beta_j(f)=0$, the scalar at each later time is the negative of the
remaining tail and has magnitude at most $C d^{-2}\|k_n\|_\infty$.
Hence the whole vector contracts at rate $(C d^{-3})^{n-j}$.  Enlarge $d$
so that this is bounded by $d^{-w(n-j)}$.  If $\beta_j(f)\ne0$, division
by $M_{j,n}$ shows that its scalar converges to the nonzero value
$\beta_j(f)$ while its kernel part tends to zero.  It therefore enters
\eqref{pure:cone}.  Finally, for a vector in that cone,
\[
 \frac{\|P_{ds}f\|_\infty}{|\ell_{ds}(P_{ds}f)|}
 \leq \frac{C+C d^{-2}M}{1-C d^{-2}M}.
\]
Choose $M$ first, larger than the limiting bound, and then $d$ large.
This proves invariance of the cone.
\end{proof}

\begin{lemma}[Energy calibration]
\label{pure:calibration}
For all sufficiently large integer $s$,
\begin{equation}
 c/s\leq\|P_sB\|_\infty\leq C/s.
 \label{pure:energy-norm}
\end{equation}
Along the geometric scales above, $P_sB$ eventually belongs to the cone
\eqref{pure:cone}.
\end{lemma}

\begin{proof}
If a transferred $B$ belonged to the stable hyperplane, then
\eqref{pure:stable-decay} and the Markov property, tested against a distant
$B_x$, would give $\E_0BB_x=O(|x|^{-w})$.  This contradicts
\eqref{pure:energy-input}, since $w>2$.  The cone assertion follows from
Lemma~\ref{pure:hyperplane}.

We next show that, for one sufficiently large fixed $T$, any centered
unit-norm cone observable $f_s$ satisfies
\begin{equation}
 \E_0[f_s\,\tau_{Ts e_1}f_s]\geq c_T>0,
 \label{pure:cone-pairing}
\end{equation}
uniformly for all sufficiently large $s$; choose $T$ divisible by $8$, and
$\tau_x$ translates an observable by $x$.  The value $c_T$ can be chosen
of order $T^{-2}$.

To see this, choose nonzero nonnegative smooth bumps $\psi_1,\psi_2$
in two disjoint regions strictly inside the unit square.  Let $Y_s$ be
the quarter-turn average of $\Phi_s(\psi_1)\Phi_s(\psi_2)$, centered
under $\E_0$.  The corresponding continuum spin-smear product,
averaged under quarter turns and centered, has nonzero scalar coefficient.
For two unsmeared pairs $z_1,z_2$ and $Te_1+y_1,Te_1+y_2$, the squared
four-spin formula of the companion gives
\[
 \Cov\bigl(\sigma(z_1)\sigma(z_2),
       \sigma(Te_1+y_1)\sigma(Te_1+y_2)\bigr)
 =\frac{|z_1-z_2|^{3/4}|y_1-y_2|^{3/4}}{4T^2}
     +o(T^{-2}).
\]
Here the neutral assignments within each pair give the product term;
their cross logarithms are $O(T^{-2})$ and cancel to first order after
sign summation.  The two assignments with charge $2$ and $-2$ on the
respective pairs give the displayed strictly positive coefficient after
taking the positive square root.  The expansion is uniform on the
separated bump supports, so it survives smearing and averaging.

Take a sufficiently large fixed $A$ and put
$V_s=\chi_A(Y_s)-\E_0\chi_A(Y_s)$, where
$\chi_A(t)=\max\{-A,\min\{A,t\}\}$.
The nonzero scalar coefficient survives for all sufficiently large $A$:
the moment bounds give ordinary $L^2$ convergence of the truncated
continuum test, and the radial coefficient is continuous in that norm.
Its measured lattice coefficient is then bounded away from zero once $H$
and the lattice threshold are large.

To see why $A$ may be fixed before $T$, put
$\varepsilon_A=\sup_s\E_0|Y_s-V_s|$, which tends to zero by the moment
bounds.  Buffered comparison gives
\[
 \|P_{2s}(Y_s-V_s)\|_\infty\leq K\varepsilon_A,\qquad
 \|P_{2s}Y_s\|_\infty+\|P_{2s}V_s\|_\infty\leq C.
\]
The last constant is independent of $A$, since
$\E_0|V_s|\leq2\E_0|Y_s|$.  Apply the overall $C/T$ transfer bound of
Lemma~\ref{pure:fixed-ratio} to both factors, now conditioned on these
radius-$2s$ collars, and integrate in the two disjoint radius-$Ts/4$
squares.  Expanding the difference of the two products yields
\[
 \bigl|\E_0[Y_s\tau_{Ts e_1}Y_s]
        -\E_0[V_s\tau_{Ts e_1}V_s]\bigr|
       \leq CT^{-2}\varepsilon_A.
\]
Thus a fixed large $A$ preserves the positive $T^{-2}$ coefficient.
For all large fixed $T$, the covariance of $V_s$ is at least $cT^{-2}$
once the lattice threshold is large enough.

Now put $a_s=\ell_s(f_s)/\ell_s(V_s)$.  The cone condition and the uniform
coefficient bounds give $0<c\leq|a_s|\leq C$.  The observable
$f_s-a_sV_s$ has measured coefficient zero.  Transfer it and the other
factor into the disjoint squares of radius $Ts/4$ about their centers.
Lemma~\ref{pure:fixed-ratio} bounds the first factor by $CT^{-3}$ and
the other centered factor by $CT^{-1}$.  Thus
\[
 \E_0[f_s\,\tau_{Ts e_1}f_s]
   =a_s^2\E_0[V_s\,\tau_{Ts e_1}V_s]+O(T^{-4}),
\]
which proves \eqref{pure:cone-pairing}.  In this argument choose the
truncation, then $T$, then a common measurement distance for the finite
list $d,T/4,T/8$, and finally the lattice threshold.  The ratio $T/8$
is used after transferring a truncation error to its radius-$2s$ collar.
Reapply Lemma~\ref{pure:hyperplane} with these choices; its constants are
uniform and its conclusions are unchanged.  For this final choice, the
covariance contradiction at the start of the proof again shows that
$P_sB$ is not stable, so it eventually enters the newly constructed cone.

Apply \eqref{pure:cone-pairing} to
$f_s=P_sB/\|P_sB\|_\infty$.  By simultaneous Markov integration in the
two disjoint squares,
\[
 \E_0[(P_sB)\,\tau_{Ts e_1}(P_sB)]
       =\E_0[B B_{Ts e_1}].
\]
The cone lower bound and \eqref{pure:energy-input} give the upper bound
in \eqref{pure:energy-norm}.  The inequality
$|\E_0[f_s\tau_{Ts e_1}f_s]|\leq1$ gives its lower bound.  We first
obtain these estimates along the geometric scales.  For an intermediate
radius, the tower property bounds its norm between the norms at the two
adjacent scales, proving \eqref{pure:energy-norm} for every large $s$.
\end{proof}

\begin{theorem}[Exact pure energy charge]
\label{pure:charge}
For every bounded even local observable $f$, the limit
\begin{equation}
 Q(f)=\kappa^{-1}\lim_{|x|\to\infty}|x|^2\E_0[fB_x]
 \label{pure:charge-definition}
\end{equation}
exists.  It defines a linear functional that kills constants, is invariant
under lattice translations and square symmetries, and satisfies $Q(b_e)=1$.
Pure transfer preserves $Q$ exactly.  If $f$ is supported in $B_s(x)$,
then
\begin{equation}
 |Q(f)|\leq Cs\|f\|_\infty.
 \label{pure:charge-bound}
\end{equation}
For centered $f$ and $R\geq4s$,
\begin{equation}
 \|P_{R,x}f\|_\infty\leq C(s/R)\|f\|_\infty.
 \label{pure:even-transfer}
\end{equation}
If $Q(f)=0$, the factor $s/R$ improves to $(s/R)^u$; it improves to
$(s/R)^w$ if, after auxiliary integration, $f$ is also invariant under
quarter turns about $x$.  More generally, with $v=u$ or $v=w$ according
to these two symmetry classes,
\begin{equation}
 \|P_{R,x}f-Q(f)P_{R,x}B_x\|_\infty
      \leq C(s/R)^v\|f\|_\infty.
 \label{pure:energy-remainder}
\end{equation}
All constants are independent of $s,R$, and $f$.
\end{theorem}

\begin{proof}
We first work at the origin on the geometric scales and with
quarter-turn-invariant centered functions.  By
Lemma~\ref{pure:calibration}, the vector $E_j=P_{s_j}B$ lies eventually
in the cone and has norm comparable to $s_j^{-1}$.  Since
$\beta_j=\ell_{s_j}+O(d^{-2})$ and $d$ is large,
\[
 c/s_j\leq|\beta_j(E_j)|\leq C/s_j.
\]
Define temporarily
$\widetilde Q(f)=\beta_j(f)/\beta_j(E_j)$ for a function supported in
$B_{s_j}$.  Exact compatibility of both numerator and denominator makes
this definition invariant under transfer to later geometric scales.
This is also consistent when the same function is regarded as belonging
to a later support space: the difference between it and its transfer
vanishes after the next transfer, hence has zero $\beta$ by compatibility.
Its kernel is exactly the stable hyperplane, and
$|\widetilde Q(f)|\leq Cs_j\|f\|_\infty$.

Increase the fixed starting scale once so that the calibration holds on
the whole grid.  For a function supported in radius $s$ above this
threshold, use the next containing geometric scale $r$, so that
$s\leq r\leq ds$.  At that scale compare $f$ with
$\widetilde Q(f)P_rB$, whose norm is at most $C\|f\|_\infty$ by
\eqref{pure:energy-norm}.  Their difference is stable.
Lemma~\ref{pure:hyperplane}, followed by norm contraction between adjacent
scales, proves \eqref{pure:energy-remainder} with $v=w$ and
$\widetilde Q$ in place of $Q$.  Radii $R$ within a fixed factor of $s$
are absorbed by increasing the constant.

For a general even function, apply this construction to its quarter-turn
average and set $\widetilde Q(f)=\widetilde Q(\mathcal A_s f)$.
The complementary component contracts at rate $(Cd^{-2})^j$ by
\eqref{pure:ratio-nonscalar}.  Choose $d$ large enough that this is at
most $d^{-uj}$.  This proves the general remainder estimate with $v=u$.
Together with $\|P_RB\|_\infty\leq C/R$, it proves
\eqref{pure:even-transfer}.

For fixed $f$, pair its remainder with $B_x$ and integrate both factors
in disjoint squares of radii proportional to $|x|$.  The remainder has
norm $O_f(|x|^{-u})$, and the transferred $B_x$ has norm $O(|x|^{-1})$.
Since $u>1$, their covariance is $o(|x|^{-2})$.
Equation~\eqref{pure:energy-input} therefore gives
\[
 |x|^2\E_0[fB_x]\longrightarrow\kappa\widetilde Q(f).
\]
This proves existence of \eqref{pure:charge-definition} and identifies
$Q=\widetilde Q$.  In particular the definition is independent of the
geometric grid and its starting scale.  Translation and square invariance
follow from this intrinsic limit and the corresponding symmetries of
$\mu_0$; the limit is unchanged by a fixed displacement of its distant
test.  Equation~\eqref{pure:energy-input} gives $Q(b_e)=1$ for every edge.
Finally the Markov property preserves the covariance with every sufficiently
distant $B_x$, so it preserves $Q$ under transfer through any containing
square.  This also establishes the stated results about arbitrary centers.
For the finitely many radii below the fixed threshold, transfer first to
that threshold and enlarge the universal constants.  Integrating independent
auxiliary variables at the outset proves the asserted extension to them.
\end{proof}

\subsection{Separated products and positive local weights}

We now use the exact normalization to control the energy component of a
product.  Write $q_s(f)=Q(f)/s$.  A product of two centered energy
observables has a much smaller charge than either factor; duality supplies
the cancellation responsible for this gain.

\begin{proposition}[Covariance and charge of separated products]
\label{pure:fusion}
Fix $D<\infty$.  Let $f,g$ be centered even observables of norm at most
one, supported respectively in $B_{Ds}(x)$ and $B_{Ds}(y)$, where
$s\geq1$ and $|x-y|=ls$.  For sufficiently large $l$, uniformly in $s,x,y$
and in the observables,
\begin{align}
 \E_0fg&=l^{-2}\bigl(\kappa q_s(f)q_s(g)+o_D(1)\bigr),
       \label{pure:fusion-covariance}\\
 |q_s(fg)|&\leq C_Dl^{-u}.
       \label{pure:fusion-charge}
\end{align}
If each observable is invariant under quarter turns about its own center
after auxiliary integration, the last bound improves to $C_Dl^{-w}$.
The term $o_D(1)$ tends to zero as $l\to\infty$ independently of $s$.
\end{proposition}

\begin{proof}
Put $r=|x-y|$ and integrate in disjoint squares of radius a small fixed
fraction of $r$.  Write $F,G$ for the transferred observables and
$U_x,U_y$ for the corresponding transfers of $B_x,B_y$.
Theorem~\ref{pure:charge} gives
\[
 \|F\|_\infty+\|G\|_\infty\leq C_Dl^{-1},\qquad
 \|F-Q(f)U_x\|_\infty+\|G-Q(g)U_y\|_\infty
       \leq C_Dl^{-v},
\]
where $v=u$, or $v=w$ in the symmetric case.  Notice that these estimates
compare with the bond observable after transfer to a containing scale;
the untransferred $Q(f)B_x$ can have norm of order $s$.
It follows that
\begin{equation}
 \bigl|\E_0fg-Q(f)Q(g)\E_0B_xB_y\bigr|
       \leq C_Dl^{-v-1}.
 \label{pure:product-covariance-error}
\end{equation}
Since $v>1$, \eqref{pure:energy-input} proves
\eqref{pure:fusion-covariance}.  Its error is uniform: $r=ls\geq l$, the
energy remainder tends uniformly to zero outside growing disks, and
$|q_s(f)|+|q_s(g)|\leq C_D$.

Simultaneous integration in the disjoint squares also preserves the
charge of their product.  Indeed it preserves its correlation with every
sufficiently distant energy test, so \eqref{pure:charge-definition}
applies.  All transferred products fit in a square of radius $Cr$.
The charge bound, applied to the same product difference as in
\eqref{pure:product-covariance-error}, therefore gives
\begin{equation}
 \frac1s\bigl|Q(fg)-Q(f)Q(g)Q(B_xB_y)\bigr|
      \leq C_D\frac rs l^{-v-1}=C_Dl^{-v}.
 \label{pure:product-charge-error}
\end{equation}

It remains to prove
\begin{equation}
 |Q(B_xB_y)|\leq Cr^{-w}.
 \label{pure:bond-fusion}
\end{equation}
For distinct edges, Kramers--Wannier duality at the critical coupling gives
\begin{equation}
 \E_0\prod_{i=1}^n b_{e_i}
        =(-1)^n\E_0^*\prod_{i=1}^n b_{e_i^*}.
 \label{pure:duality-centered}
\end{equation}
Here $e_i^*$ is the crossing dual edge.  To check the centering directly,
differentiate the finite planar partition-function duality
\cite{KW} in the distinct couplings.  Its coupling-dependent prefactor is
$\prod_e(\sinh(2K_e))^{1/2}$, and
$\sinh(2K_e)\sinh(2K_e^*)=1$.  At $K_0$ the logarithmic prefactor
derivative is $\sqrt2$ and $\dd K_e^*/\dd K_e=-1$; thus
$\partial_{K_e}-1/\sqrt2$ becomes
$-(\partial_{K_e^*}-1/\sqrt2)$.  Distinct derivatives introduce no contact
terms.  Exhausting the plane and using uniqueness of its pure critical
state proves \eqref{pure:duality-centered}.

Apply \eqref{pure:duality-centered} to the three separated bond averages
in the definition of $Q(B_xB_y)$.  After identifying the dual lattice
with the original lattice, the distant dual average $B'_z$ is centered
and has charge one.  Thus $B'_z-B_z$ is centered and has charge zero.
With $x,y$ fixed, separately transfer the centered product
$B'_xB'_y-\E_0(B'_xB'_y)$ and this distant difference into squares of
radii proportional to $|z|$.  Their norms are respectively
$O_{x,y}(|z|^{-1})$ and $O(|z|^{-u})$, by
Theorem~\ref{pure:charge}.  Their covariance is therefore
$O_{x,y}(|z|^{-u-1})=o(|z|^{-2})$.  Replacing the distant dual average
by $B_z$ leaves the charge limit unchanged, and consequently
\[
 Q(B_xB_y)=-Q(B'_xB'_y),
\]
where the primed average consists of the four dual edges surrounding the
original site.

There are four identifications obtained by the translations
$h=(\pm\frac12,\pm\frac12)$.  They send the dual average about $x$
to the boundary-edge averages of the four unit plaquettes incident to $x$.
Let $B'_{i,x}$, $1\leq i\leq4$, denote the resulting observables, and put
\[
 C_x=\frac14\sum_{i=1}^4B'_{i,x},\qquad D_{i,x}=B'_{i,x}-C_x.
\]
The observable $C_x$ is quarter-turn invariant about $x$, has charge one,
and is supported in a fixed-radius square.  Each $D_{i,x}$ has charge zero,
and their average vanishes.  Hence
\[
 \frac14\sum_iQ(B'_{i,x}B'_{i,y})
   =Q(C_xC_y)+\frac14\sum_iQ(D_{i,x}D_{i,y}).
\]
The charge bound after separate transfers gives
\[
 |Q(D_{i,x}D_{i,y})|\leq Cr^{1-2u},\qquad
 |Q(C_xC_y)-Q(B_xB_y)|\leq Cr^{-w}.
\]
For the first estimate both factors have zero charge and decay as $r^{-u}$.
For the second, expand the difference with one factor $C-B$; it is
quarter-turn invariant and has zero charge, while the other centered
factor decays as $r^{-1}$.  The common-support charge bound supplies the
factor $Cr$ in each case.  Averaging the negative duality identity now
gives $2|Q(B_xB_y)|\leq C(r^{-w}+r^{1-2u})$.  Since $2u-1>w$, this is
\eqref{pure:bond-fusion}.

Finally, its contribution in \eqref{pure:product-charge-error} is at most
\[
 \frac{|Q(f)Q(g)|}{s}\,Cr^{-w}
       \leq C_Ds(ls)^{-w}
       =C_Ds^{1-w}l^{-w}\leq C_Dl^{-w}.
\]
This proves both charge bounds uniformly for $s\geq1$.
\end{proof}

\begin{lemma}[Transfer of a positive local weight]
\label{pure:positive-weight}
There is a constant \(C<\infty\) with the following property.
Let \(m,R\) be integers with \(m\ge1\) and \(R\ge8m\).
Let \(\eta\) be a finite family of independent auxiliary variables,
with arbitrary reference laws, independent of the full-plane pure
critical Ising field. Suppose that \(G=G(\sigma_{B_m},\eta)\) is a bounded measurable
real function satisfying \(G(-\sigma,\eta)=G(\sigma,\eta)\).
Write
\[
 \|G\|_\infty
 =\max_{\sigma_{B_m}}\operatorname*{ess\,sup}_{\eta}
       |G(\sigma_{B_m},\eta)|,
 \qquad
 \mathcal Z_R^G(\tau)
 =\E_0\!\left[e^G\mid\sigma_{\partial B_R}=\tau\right],
\]
where \(\E_0\) includes integration over all the auxiliary variables
used by \(G\). Then
\begin{equation}
 \osc_{\tau}\log\mathcal Z_R^G(\tau)
 \le C\frac mR\min\{1,\|G\|_\infty\}.
 \label{pure:positive-weight-bound}
\end{equation}
Here \(\osc\) is the supremum minus the infimum over all boundary
spin assignments, and \(C\) is independent of the number and laws of
the auxiliary variables. For a patch of real radius \(a\ge0\), one
may take \(m=\max\{1,\lceil a\rceil\}\).
\end{lemma}

\begin{proof}
First integrate the auxiliary variables and put
\[
 W(\sigma_{B_m})=\E_{\eta}e^{G(\sigma_{B_m},\eta)},
 \qquad Z=\E_0W,
 \qquad
 h(\xi)=\frac{\E_0[W\mid\sigma_{\partial B_{2m}}=\xi]}{Z}.
\]
The function \(W\) is positive and even. The fixed-ratio buffered
comparison of \cite[Section~2, Proposition ``Buffered and small-patch
comparisons'']{S}, applied to \(B_m\) and \(\partial B_{2m}\), gives
a constant \(K\ge1\) such that
\[
 K^{-1}\le h(\xi)\le K
 \quad\text{for every }\xi.
\]
Indeed, the conditional law on \(B_m\) has pointwise density between
\(K^{-1}\) and \(K\) relative to its full-plane marginal: the latter
is a mixture of the same conditional laws. Integrating the positive
function \(W\) proves the displayed bounds, regardless of its size.

Spin-flip symmetry gives \(h(-\xi)=h(\xi)\), and the tower property
gives \(\E_0h=1\). If \(t=\|G\|_\infty\), then both \(W\) and
\(Z\) lie between \(e^{-t}\) and \(e^t\). Consequently
\[
 \|h-1\|_\infty
 \le \min\{K+1,e^{2t}-1\}
 \le C_K\min\{1,t\}.
\]
Apply the even transfer estimate in Theorem~\ref{pure:charge} to the
centered even function \(h-1\), supported in \(B_{2m}\). Since
\(R\ge8m\), it gives
\[
 \|P_R(h-1)\|_\infty
 \le C\frac mR\min\{1,t\}.
\]
The Ising Markov property, followed by conditional expectation, yields
\[
 \frac{\mathcal Z_R^G(\tau)}{Z}
 =\E_0[h\mid\sigma_{\partial B_R}=\tau]
 =1+P_R(h-1)(\tau)\ge K^{-1}.
\]
The logarithm is \(K\)-Lipschitz on \([K^{-1},\infty)\).
Its oscillation in the last display is therefore at most
\(2K\|P_R(h-1)\|_\infty\), which proves
\eqref{pure:positive-weight-bound}.
\end{proof}
